%%=============================================================================
%% CAPITULO 5 -- CONSIDERACOES FINAIS
%%=============================================================================

\section{Considerações finais}

Faça uma reflexão sobre a questão problema. Verifique se sua hipótese e seus
objetivos geral e específicos foram alcançados. Apresente suas sugestões e
cuidados, caso outras pessoas queiram implantar as ações em suas atividades
profissionais.

As considerações finais não devem introduzir dados ou argumentos novos: elas
retomam o percurso da pesquisa, respondem à pergunta formulada na introdução,
reconhecem as limitações do estudo e apontam possibilidades de investigações
futuras.

A formatação deste trabalho segue as normas da ABNT para apresentação de
trabalhos acadêmicos \parencite{abnt14724}, para elaboração de referências
\parencite{abnt6023} e para apresentação do sumário \parencite{abnt6027}. Note,
na lista de referências, que a autoria é repetida por extenso, conforme os
exemplos do template institucional Word do IDP de 2026, que prevalece sobre
o manual nos pontos de divergência.
